% Automatisch aus src/evaluation/indicator_map.py (MAPPING) generiert.
% Bei Aenderungen an der Zuordnung: Skript erneut laufen lassen, nicht von Hand nachpflegen.
\begingroup
\tiny
\renewcommand{\arraystretch}{1.15}

\setlength{\LTleft}{-1.5cm}
\setlength{\LTright}{0pt}

\begin{longtable}{p{3.1cm} p{0.9cm} p{3.6cm} p{6.3cm}}
\label{tab:mqa_abc_mapping}\\

\toprule
\textbf{Indikator} & \textbf{Klasse} & \textbf{MQA-Metrik(en)} & \textbf{Begründung} \\
\midrule
\endfirsthead

\toprule
\textbf{Indikator} & \textbf{Klasse} & \textbf{MQA-Metrik(en)} & \textbf{Begründung} \\
\midrule
\endhead

\midrule
\multicolumn{4}{r}{\textit{Fortsetzung auf der nächsten Seite}} \\
\midrule
\endfoot

\bottomrule
\endlastfoot

\multicolumn{4}{l}{\textbf{Auffindbarkeit}} \\
\midrule

find\_keywords\_count & B\,\textcolor{red}{$\boldsymbol{\ast}$} & keyword\_availability & MQA: Presence dcat:keyword; Prototyp: Anzahl im Band 3$\leq$k$\leq$15 (ternär, Handreichung). \\

find\_theme\_valid & B & category\_availability & MQA: Presence dcat:theme; Prototyp: Wert aus kontrolliertem Vokabular. \\

find\_temporal\_coverage & A\,\textcolor{red}{$\boldsymbol{\ast}$} & temporal\_availability & Beide: zeitliche Abdeckung (start/end). Prototyp prüft zusätzlich Datentyp. \\

find\_locn\_geometry & B & spatial\_availability & MQA: dct:spatial Presence; Prototyp: tatsächliche locn:geometry (nicht leer). \\

find\_political\_geocoding & B & spatial\_availability & MQA: dct:spatial Presence; Prototyp: dcatde:politicalGeocodingURI aus Vokabular (Fallback locn:adminUnitL2 $\rightarrow$ nur PARTIAL). \\

find\_geocoding\_level & C & --- & MQA: keine Metrik für die administrative Ebene (dcatde:politicalGeocodingLevelURI, Konvention 09). \\

find\_accrual\_periodicity$^{\dagger}$ & C & --- & MQA: keine Metrik für Aktualisierungsfrequenz. \\

find\_issued\_datetime & A\,\textcolor{red}{$\boldsymbol{\ast}$} & date\_issued\_availability & Beide: Vorhandensein dct:issued. Prototyp prüft zusätzlich xs:date/dateTime-Typ. \\

find\_modified\_datetime & A\,\textcolor{red}{$\boldsymbol{\ast}$} & date\_modified\_availability & Beide: Vorhandensein dct:modified. Prototyp prüft zusätzlich xs:date/dateTime-Typ. \\


\multicolumn{4}{l}{\textbf{Zugänglichkeit}} \\
\midrule

acc\_download\_url & B\,\textcolor{red}{$\boldsymbol{\ast}$} & download\_url\_availability & MQA: Presence dcat:downloadURL ($\geq$1 Distribution). Prototyp: Anteil ALLER Distributionen mit downloadURL (graded) — strenger, PASS erst wenn nahezu alle eine haben. \\

acc\_download\_url\_response & A & download\_url\_status\_code & Beide: HTTP-Auflösbarkeit der downloadURL ($<$400). \\

acc\_access\_url\_response & A & access\_url\_status\_code & Beide: HTTP-Auflösbarkeit der accessURL ($<$400). \\

acc\_format & A\,\textcolor{red}{$\boldsymbol{\ast}$} & format\_availability \textsc{und} format\_media\_type\_vocabulary & MQA: Format vorhanden + format\_media\_type\_vocabulary (Format OR MediaType im Vokabular, kombiniert).  \\

acc\_media\_type & B\,\textcolor{red}{$\boldsymbol{\ast}$} & media\_type\_availability \textsc{und} format\_media\_type\_vocabulary & MQA: MediaType vorhanden + format\_media\_type\_vocabulary (Format OR MediaType im Vokabular, kombiniert). Prototyp: Anteil der Distributionen mit MediaType als IANA-URI im IANA-Vokabular (graded) — strenger, da MQA keine IANA-URI-Form prüft und der Vocab-Check via Format-Treffer auch ohne MediaType-Eintrag PASS gibt. \\

acc\_machine\_readable\_access & B & format\_machine\_readable & MQA: Listen-Lookup beste Distribution; Prototyp: Mittelwert aller Distributionen (high +1 / mid +0.5 / none 0, kein Malus). \\

acc\_format\_non\_proprietary & A & format\_non\_proprietary & Beide: nicht-proprietäres Format-URI aus EU-Vokabular. Prototyp: Anteil aller Distributionen. \\


\multicolumn{4}{l}{\textbf{Wiederverwendbarkeit}} \\
\midrule

reuse\_dcat\_ap\_de\_compliance & A & dcat\_ap\_compliance & Beide: DCAT-AP SHACL-Konformität via ITB-API (wenn MQA shacl\_validation=True). FAIL bei mindestens einer Verletzung. \\

reuse\_license & B\,\textcolor{red}{$\boldsymbol{\ast}$} & license\_availability, known\_license & MQA: Lizenz vorhanden (license\_availability) + EU-Autoritäts-URI (known\_license). Prototyp: je Distribution nur freie DCAT-AP-DE-Lizenz = +1, eingeschränkt/unbekannt/fehlt = -0.5 Malus, dann Mittel (graded, distributionsebene) $\rightarrow$ strenger. \\

reuse\_access\_rights$^{\dagger}$ & A\,\textcolor{red}{$\boldsymbol{\ast}$} & access\_rights\_availability \textsc{und} access\_rights\_vocabulary & Beide: accessRights vorhanden + Vokabular. \\

reuse\_publisher & A & publisher\_availability & Beide prüfen auf Präsenz des dct:publisher \\

reuse\_contact & B & contact\_point\_availability & MQA: Presence dcat:contactPoint; Prototyp: valide vcard:hasEmail ODER vcard:hasURL (Konvention 01), sonst PARTIAL. \\

reuse\_availability & C & --- & MQA: keine Verfügbarkeits-Metrik; Prototyp: Anteil der Distributionen mit dcatap:availability aus dem Planned-Availability-Vokabular. \\

reuse\_contributor\_id & C & --- & MQA: keine Metrik (DCAT-AP-DE dcatde:contributorID). \\


\multicolumn{4}{l}{\textbf{Aussagekraft}} \\
\midrule

expr\_title\_quality & C & --- & MQA: keine Titel-Inhaltsprüfung. \\

expr\_description\_quality & C & --- & MQA: keine Beschreibungs-Inhaltsprüfung. \\

expr\_keyword\_quality & C & --- & MQA: zählt Keywords, prüft keine Bedeutung. \\

expr\_title\_description\_coherence & C & --- & MQA: keine Kohärenzprüfung Titel$\leftrightarrow$Beschreibung. \\

expr\_thematic\_consistency & C & --- & MQA: keine thematische Konsistenzprüfung. \\

expr\_contextual\_qualifiers & C & --- & MQA: keine Prüfung kontextueller Qualifizierer. \\


\end{longtable}
\endgroup

{\footnotesize $^{\ast}$ Grenzfall, gegen die MQA-Spezifikation zu pruefen (\texttt{review=True} in \texttt{indicator\_map.py}). $^{\dagger}$ In der finalen Konfiguration \texttt{state\_evaluation\_final.yaml} per Indikator-Blacklist deaktiviert; taucht in den Kapitel-5-Kennzahlen nicht auf (dort 27 statt 29 klassifizierte Indikatoren). \par}